\documentclass{article}
\usepackage{amsmath, amssymb, amsthm}
\usepackage{hyperref}
\usepackage{geometry}
\geometry{margin=1in}

\title{Erd\H{o}s Problem \#929}
\date{Last updated: 2025-08-31}
\author{Source: erdosproblems.com}

\begin{document}
\maketitle

\noindent\textbf{Status:} OPEN \quad \textbf{No prize}

\noindent\textbf{Tags:} number theory

\medskip
\noindent\textbf{Problem Statement:}

Let $k\geq 2$ be large and let $S(k)$ be the minimal $x$ such that there is a positive density set of $n$ where\[n+1,n+2,\ldots,n+k\]are all divisible by primes $\leq x$.

Estimate $S(k)$ - in particular, is it true that $S(k)\geq k^{1-o(1)}$?

\bigskip
\noindent\textbf{Remarks:}

It follows from Rosser's sieve that $S(k)> k^{1/2-o(1)}$.

It is trivial that $S(k)\leq k+1$ since, for example, one can take $n\equiv 1\pmod{(k+1)!}$. The best bound on large gaps between primes due to Ford, Green, Konyagin, Maynard, and Tao \cite{FGKMT18} (see [4]) implies\[S(k) \ll k \frac{\log\log\log k}{\log\log k\log\log\log\log k}.\]

\bigskip
\noindent\textbf{References:}

[FGKMT18] Ford, Kevin and Green, Ben and Konyagin, Sergei and Maynard, James and Tao, Terence, Long gaps between primes. J. Amer. Math. Soc. (2018), 65-105.

\bigskip
\noindent\small{Source: \url{https://www.erdosproblems.com/929}}
\end{document}
